% --- Editorial: Volume 2025.1 ---
% This file contains the letter from the editor for the inaugural issue.

\documentclass[conference, compsoc, 11pt]{IEEEtran}
\usepackage{url}

% orcid icon (TODO: move to a separate style file)
\usepackage{scalerel}
\usepackage{tikz}
\usetikzlibrary{svg.path}

\definecolor{orcidlogocol}{HTML}{A6CE39}
\tikzset{
  orcidlogo/.pic={
    \fill[orcidlogocol] svg{M256,128c0,70.7-57.3,128-128,128C57.3,256,0,198.7,0,128C0,57.3,57.3,0,128,0C198.7,0,256,57.3,256,128z};
    \fill[white] svg{M86.3,186.2H70.9V79.1h15.4v48.4V186.2z}
                 svg{M108.9,79.1h41.6c39.6,0,57,28.3,57,53.6c0,27.5-21.5,53.6-56.8,53.6h-41.8V79.1z M124.3,172.4h24.5c34.9,0,42.9-26.5,42.9-39.7c0-21.5-13.7-39.7-43.7-39.7h-23.7V172.4z}
                 svg{M88.7,56.8c0,5.5-4.5,10.1-10.1,10.1c-5.6,0-10.1-4.6-10.1-10.1c0-5.6,4.5-10.1,10.1-10.1C84.2,46.7,88.7,51.3,88.7,56.8z};
  }
}

\newcommand\orcidicon[1]{\href{https://orcid.org/#1}{\mbox{\scalerel*{
\begin{tikzpicture}[yscale=-1,transform shape]
\pic{orcidlogo};
\end{tikzpicture}
}{|}}}}

\usepackage{hyperref} %<--- Load after everything else

% --- Document Metadata ---
\title{Editorial: Building a New Commons for Scholarly Work}

\author{
    \IEEEauthorblockN{Ivan Anishchuk}
    \IEEEauthorblockA{
        ivan@ivananishchuk.net https://ivananishchuk.work \\
	ORCID \orcidicon{0009-0008-8729-6137}: 0009-0008-8729-6137 % Placeholder
    }
}

\begin{document}

% FIXME: provide correct publication info when available
\IEEEoverridecommandlockouts
\IEEEpubid{\makebox[\columnwidth]{25.1.araxiv.org/25.1.2~CC0 2025 Ivan Anishchuk \hfill}
\hspace{\columnsep}\makebox[\columnwidth]{ }}

\maketitle

\IEEEpubidadjcol

% --- Abstract ---
\begin{abstract}
An introduction to the inaugural release of the Transactions of the Amateur Research Association. This editorial frames the project's mission within the context of the high costs of academic publishing and the rise of the independent researcher. It outlines the ARA's philosophy of leveraging open-source tools and automation to build a new, accessible, and legitimate venue for scholarly communication, and presents the foundational whitepaper as the issue's centerpiece.
\end{abstract}

% --- Keywords ---
\begin{IEEEkeywords}
editorial, decentralized science, open science, scholarly communication, citizen science, open infrastructure, public goods
\end{IEEEkeywords}

% --- Main Body ---
\section*{Editorial: Building a New Commons for Scholarly Work}

The landscape of knowledge creation is experiencing a profound paradox. On one hand, the tools for research, analysis, and collaboration have never been more accessible. On the other, the formal channels for validating and disseminating this knowledge remain largely constrained by a legacy system characterized by high costs, long delays, and institutional barriers. The scholarly publishing market, valued in the tens of billions of dollars, has developed an economic model where access—both to read and to publish—is often prohibitively expensive, with average Article Processing Charges (APCs) frequently exceeding thousands of dollars \cite{solomon2016}. This creates a significant barrier for a vast and growing community of researchers who operate outside the traditional, well-funded university system.

This paradox has given rise to a new class of contributor: the unaffiliated scholar. This diverse group includes industry professionals whose work has scholarly merit, citizen scientists contributing to massive data collection and analysis efforts, and independent researchers pursuing knowledge outside of formal employment. The value generated by this community is not trivial. For instance, citizen science initiatives now contribute to a substantial volume of published research, often providing data at a scale that would be impossible for conventional research groups to gather \cite{theobald2015}. Yet, this work is frequently relegated to the gray literature of preprints, personal blogs, or internal reports, lacking the formal markers of legitimacy that would ensure its long-term discovery and integration into the scholarly record.

The primary barrier for these researchers is not a lack of quality, but a lack of access to the \textit{infrastructure of legitimacy}. This infrastructure includes services we often take for granted: persistent Digital Object Identifiers (DOIs) from agencies like Crossref, professional typesetting, inclusion in curated indexes, and a clear, verifiable institutional affiliation. The movement for "open scholarly infrastructure" argues that these core services should be community-governed and financially accessible, forming a global "research commons" that serves knowledge creation rather than commercial extraction \cite{bilder2015}. The Amateur Research Association (ARA) is a direct, practical application of this principle.

The ARA is our proposed solution: a framework for a publishing platform that is as legitimate as a formal journal but as accessible as a preprint server. Our model is built on a foundation of open-source tools and transparent, automated workflows. By leveraging ubiquitous technologies like Git for submission and version control, we embrace a culture of reproducibility and transparency that is increasingly seen as essential for rigorous scientific work \cite{noble2009}. By standardizing on LaTeX, we ensure that every publication meets a high standard of professional typesetting. Our automation pipeline, running on GitHub Actions, handles the technical burdens of validation, compilation, and deployment, allowing our human editors to focus on the intellectual curation of each release.

In this inaugural release, we are proud to present our foundational whitepaper. It is not merely a description of the ARA but its first formal artifact—a detailed blueprint for the technology, governance, and philosophy that will guide us. It outlines a phased, sustainable approach to growth, from a minimal viable product to a professionally managed, community-governed institution.

We stand at a unique moment where the tools of decentralization—both technical and social—can be leveraged to build new, more equitable institutions. The ARA is an experiment to see if a small, dedicated community can create a self-sustaining scholarly commons. We invite you to read our plan, challenge our assumptions, join our community, and, most importantly, submit your work. Together, we can build a new venue for the vital research being done at the frontiers of knowledge, regardless of institutional affiliation.


% --- References ---
% For a self-contained document, we use the thebibliography environment.
\begin{thebibliography}{1}

\bibitem{solomon2016}
D. J. Solomon and B. Bj{\"o}rk, "Peer review in a world of open access," \textit{Journal of the Association for Information Science and Technology}, vol. 67, no. 1, pp. 5--15, 2016.

\bibitem{theobald2015}
E. J. Theobald et al., "Global change and local solutions: Tapping the unrealized potential of citizen science for biodiversity research," \textit{Biological Conservation}, vol. 181, pp. 236--244, 2015.

\bibitem{bilder2015}
G. Bilder, J. Lin, and C. Neylon, "Principles for open scholarly infrastructure-v1," \textit{Figshare}, 2015. [Online]. Available: \url{https://doi.org/10.6084/m9.figshare.1314859}

\bibitem{noble2009}
W. S. Noble, "A quick guide to organizing computational biology projects," \textit{PLoS computational biology}, vol. 5, no. 7, p. e1000424, 2009.

\end{thebibliography}

\end{document}
